\section{Datasets and evaluation protocol}
\label{sec:data}
\subsection{MovieLens-100K and implicit conversion}
The initial benchmark uses the official MovieLens-100K release
\citep{harper2015movielens}. Ratings of four or five are treated as positive
implicit feedback. Ratings below four are not positive training events, but
they remain part of the observed history used for candidate filtering. No
rating-regression target is optimized. Users require at least three positives
so that one training positive and two held-out positives can remain. One user
with no positives is removed. Items are not removed for having few or no
training positives.

\begin{table}[tbp]\centering
\caption{Dataset sizes under the frozen chronological implicit-feedback protocol.}
\label{tab:data}
\begin{tabular}{lrr}\toprule
Quantity & MovieLens-100K & MovieLens-1M \\\midrule
Eligible users & 942 & 6,035 \\
Catalog items appearing in ratings & 1,682 & 3,706 \\
Training positives & 53,491 & 563,206 \\
Validation targets & 942 & 6,035 \\
Test targets & 942 & 6,035 \\
Test targets scored in this project & 942 & 0 \\
\bottomrule\end{tabular}
\tabnote{The ML-100K raw release contains 943 users and 100,000 ratings;
55,375 ratings are positive. ML-1M has 575,281 positives before eligible-user
filtering; five users are removed. The ML-1M catalog is the set represented
in the ratings file, not every row of a metadata catalog.}
\source{Frozen dataset report and E3 dataset provenance JSON.}
\end{table}

Each user's positives are ordered by timestamp. The last positive becomes
the test target, the penultimate positive the validation target, and the rest
training observations. Equal timestamps use a deterministic SHA-256 tie breaker
based on user ID, item ID, and split seed 2026. This is a per-user chronology,
not a global temporal split across all users. Ties cannot recover the actual
order of simultaneous events. On ML-100K, 316 users have a validation timestamp
equal to their final training timestamp, and 305 have equal validation and test
timestamps; the fixed convention is therefore consequential and documented.

The retained ML-100K users have a mean of 56.7845 and median of 37.5 training
positives, with a range of 1--376. Activity groups are defined solely from these
counts: low activity has at most 22 positives (319 users), medium activity
23--61 (310 users), and high activity at least 62 (313 users). Six validation
and fourteen test targets are cold items with no training-positive observation
across the population. They remain in the principal evaluation.

\subsection{Negative sampling and candidate filtering}
For training, a negative is sampled uniformly from items outside that user's
training-positive set. The sampler does not consult validation or test labels.
Consequently a future positive can occasionally be sampled as unobserved; filtering
it out would use held-out information. Low-rated items can also be sampled as
ordinary nonpositives. This convention is distinct from the evaluation filter.

For each held-out event, evaluation ranks the complete catalog after removing
all items the user rated before that event, including low-rated items. The
target is retained. Later ratings never filter an earlier event's candidate
set. Validation and test therefore have different historical filters, and
the test history includes the validation event. On ML-100K the mean candidate
counts are 1,581.5658 and 1,579.0658 respectively. There is no sampled-negative
evaluation that could make ranking artificially easy or alter comparability
between methods.

Metric ties are resolved pessimistically for the relevant target. Diagnostic
top-10 sets in E2--E4 instead use a deterministic item-index tie rule and exclude
training positives only. Those label-free diagnostic candidate sets are not
interchangeable with the historical held-out evaluator.

\subsection{Ranking metrics}
With one relevant held-out item per user and its one-based rank $R_u$,
\begin{align}
 \mathrm{HR@10}_u&=\mathbb{1}\{R_u\leq10\},\nonumber\\
 \mathrm{NDCG@10}_u&=\frac{\mathbb{1}\{R_u\leq10\}}{\log_2(R_u+1)},\qquad
 \mathrm{MRR@10}_u=\frac{\mathbb{1}\{R_u\leq10\}}{R_u}.
 \label{eq:metrics}
\end{align}
Metrics are averaged equally over eligible users. Recall@10 equals HR@10 in
this one-target design, so it does not add an independent outcome. NDCG@10 is
the declared primary measure. HR@10 and MRR@10 are retained to expose changes
in hit rate and rank position; they are secondary, without an additional
family of confirmatory significance claims.

\subsection{MovieLens-1M replication}
E3 obtains MovieLens-1M from the official GroupLens distribution and applies
the unchanged conversion and splitting functions. The archive passes its
published MD5 and ZIP integrity checks, and raw and processed SHA-256 hashes
are recorded. \Cref{tab:data} gives the resulting sizes. E3 observes a fixed
sample of 128 users for trajectory diagnostics while allowing every eligible
user to participate in training. Full-population validation is scored for
the 30 trained models; the test set is not scored.

Both datasets originate from MovieLens. Their user or catalog overlap has
not been audited. Replication on the larger dataset is evidence that the
diagnostic is not confined to one saved small-data training run; it is not
proof of independent populations or cross-domain generalization.

\subsection{Holdout exposure and dataset limitations}
The ML-100K test results have been used repeatedly across historical stages.
An earlier regression-test incident also performed nine two-round real-data
B2 test evaluations. Those accidental outputs were not inspected for selection,
but their existence means the benchmark cannot be called a pristine holdout.
Later stages freeze their own choices before scoring, preserving a meaningful
local separation of selection and evaluation without erasing this history.

MovieLens is a public research dataset rather than a deployment privacy
population. Positive-threshold conversion, per-user holdout, missing-not-at-random
exposure, and very few cold targets all constrain interpretation. The report
does not infer a reliable cold-start improvement from fourteen examples, nor
does it treat offline rank utility as user satisfaction in a live service.
